\documentclass[11pt,a4paper]{article}
\usepackage[T1]{fontenc}
\usepackage[utf8]{inputenc}
\usepackage{amsmath,amssymb,xcolor}
\usepackage[margin=25mm]{geometry}
\usepackage[colorlinks=true,urlcolor=blue]{hyperref}
\definecolor{nexorablue}{HTML}{2459B5}
\setlength{\parindent}{0pt}
\setlength{\parskip}{7pt}
\setlength{\emergencystretch}{2em}
\renewcommand{\familydefault}{\sfdefault}
\begin{document}

{\small\bfseries\color{nexorablue} NEXORA CPGE}\par
{\small MPSI / NEX-ENS-03 / Énoncé}\par
{\LARGE\bfseries Quand deux ensembles disjoints ont une image commune}\par
\vspace{0.25em}\hrule\vspace{0.6em}

Soit $f:\mathbb R\to\mathbb R$ définie par $f(x)=(x-1)^2$. Calculer $f([-1,0])$, $f([2,3])$ et $f^{-1}([1,4])$. Expliquer pourquoi l’image directe ne préserve pas ici l’intersection.


\end{document}
